\documentclass[11pt,reqno]{amsart}

\usepackage[T1]{fontenc}
\usepackage{lmodern}
\usepackage{microtype}
\usepackage{amsmath,amssymb,mathtools}
\usepackage{booktabs}
\usepackage{enumitem}
\usepackage{xcolor}
\usepackage[colorlinks=true,linkcolor=blue!55!black,citecolor=blue!55!black,urlcolor=blue!65!black]{hyperref}
\hypersetup{
  pdftitle={An Explicit Four-Point Improvement to the Unconditional Montgomery--Taylor Bound for Simple Zeta Zeros},
  pdfauthor={Jizhou Guo},
  pdfsubject={Analytic number theory and formal verification},
  pdfkeywords={Riemann zeta function, simple zeros, pair correlation, Lean}
}

\newtheorem{theorem}{Theorem}[section]
\newtheorem{proposition}[theorem]{Proposition}
\newtheorem{lemma}[theorem]{Lemma}
\newtheorem{corollary}[theorem]{Corollary}
\theoremstyle{definition}
\newtheorem{definition}[theorem]{Definition}
\theoremstyle{remark}
\newtheorem{remark}[theorem]{Remark}

\newcommand{\R}{\mathbb{R}}
\newcommand{\eps}{\varepsilon}
\newcommand{\one}{\mathbf{1}}
\newcommand{\pos}[1]{(#1)_{+}}
\newcommand{\Ntot}{N}
\newcommand{\Nsimple}{N_{0}^{\mathrm{s}}}
\DeclareMathOperator{\tr}{tr}
\DeclareMathOperator{\rank}{rank}

\title[An explicit four-point improvement]{An Explicit Four-Point Improvement to the Unconditional Montgomery--Taylor Bound for Simple Zeta Zeros}

\author{Jizhou Guo}
\address{Dots Studio, RedNote}
\email{mitsuha2021b@gmail.com}

\subjclass[2020]{11M06, 11M26, 15A42, 68V20}
\keywords{Riemann zeta function, simple zeros, critical line, pair correlation, Montgomery--Taylor kernel, formal verification}

\begin{document}

\begin{abstract}
Let $N(T,2T)$ count nontrivial zeros of the Riemann zeta function with
$T<\operatorname{Im}\rho\le 2T$, with multiplicity, and let
$N_0^{\mathrm{s}}(T,2T)$ count those zeros that are simple and lie on the
critical line.  Building on the unconditional trace-and-pair-correlation
framework of Claude, whose optimized endpoint constant is
$H_*=0.672500703679\ldots$, we prove the explicit strict improvement
\[
 \liminf_{T\to\infty}
 \frac{N_0^{\mathrm{s}}(T,2T)}{N(T,2T)}
 \ge H_*+\frac{1}{79\,828\,975}.
\]
The new ingredient is a width-six, four-point packing argument.  We localize
small values of the normalized Montgomery--Taylor kernel near its first six
positive zeros, prove an exact three-coordinate energy lower bound
\[
 \delta_{1,6}\ge
 \frac{2810232522752}{834437357315980975},
\]
and combine the double counting of six Gram edges with a sharp trace-zero
four-block inequality.  Every terminal comparison is rational.  The full
extension, including the fixed-window transfer and endpoint passage, is
checked in Lean at a pinned public commit; executable ledgers independently
enumerate all $6^3=216$ root-index budgets.  We also determine the ceiling of
this frozen finite interface: the next unit fraction $1/79\,828\,974$ cannot
be obtained without strengthening at least one of its inputs.
\end{abstract}

\maketitle

\section{Introduction}

Let $\rho=\beta+i\gamma$ run over the nontrivial zeros of the Riemann zeta
function.  We write
\[
 N(T,2T)=\#\{\rho:T<\gamma\le 2T\},
\]
counting multiplicity, and let $N_0^{\mathrm{s}}(T,2T)$ count the simple
zeros in the same range with $\beta=1/2$.

Montgomery's pair-correlation argument gives, under the Riemann hypothesis,
the celebrated lower proportion $2/3$ of simple zeros
\cite{Montgomery1973}.  Optimizing the test function yields the
Montgomery--Taylor constant $0.6725\ldots$ \cite{Montgomery1975,CheerGoldston1993}.
Subsequent conditional improvements use additional information or wider
positivity, including the semidefinite-programming bound of Chirre,
Gon\c{c}alves and de Laat \cite{ChirreGoncalvesDeLaat2020}.  Separately,
an unconditional Montgomery form-factor theorem summing over all complex
zeros was proved by Baluyot, Goldston, Suriajaya and Turnage-Butterbaugh
\cite{BaluyotEtAl2024}, and recent work has clarified what
would follow from removing the hypothesis from the zero-side argument
\cite{BaluyotEtAl2025,GoldstonSuriajaya2025,GoldstonSuriajaya2026}.

The parent manuscript \cite{Claude2026} replaces the missing zero-side
positivity by finite-dimensional Hermitian linear algebra.  It proves
unconditionally that the Montgomery--Taylor endpoint is attained for zeros
that are simultaneously simple and on the critical line.  The present paper
is a companion improvement to that theorem.  We do not alter its prime-side
or explicit-formula analysis.  Instead, we strengthen the local spectral
defect that is inserted into the same analytic transfer.

Define
\begin{equation}\label{eq:cstar}
 c_*=
 \frac{\sqrt2\tan(1/\sqrt2)}{1+\tan(1/\sqrt2)/\sqrt2},
 \qquad H_*=2-\frac1{c_*}.
\end{equation}
Thus $H_*=0.672500703679411\ldots$.

\begin{theorem}[Main theorem]\label{thm:main}
For every $\eta>0$ there is $T_0$ such that, for all $T\ge T_0$,
\begin{equation}\label{eq:main-epsilon}
 \left(H_*+\frac1{79\,828\,975}-\eta\right)N(T,2T)
 \le N_0^{\mathrm{s}}(T,2T).
\end{equation}
Consequently,
\begin{equation}\label{eq:main-liminf}
 \liminf_{T\to\infty}
 \frac{N_0^{\mathrm{s}}(T,2T)}{N(T,2T)}
 \ge H_*+\frac1{79\,828\,975}.
\end{equation}
\end{theorem}

Numerically, the certified increment is
$1.2526779906669\cdot10^{-8}$.  Its small size should not obscure the
structural point: the two-trace endpoint is not rigid.  A local energy defect
survives all off-line Hermitian cancellations and gives a strictly positive
unconditional gain.

The proof has three new components.
\begin{enumerate}[label=(\roman*)]
\item We localize every sufficiently small value of the endpoint kernel on
      $[0,12\pi]$ inside radius $1/8$ of one of its first six positive zeros.
\item For three coordinates $a,b,a+b$, exact separation of the localized
      roots gives a uniform rational energy lower bound.
\item We pack four simple-zero labels at a time.  The four constituent
      triples count the six Gram edges twice, while a trace-zero Hermitian
      four-block converts Frobenius energy to trace-norm mass with factor two.
\end{enumerate}

All transcendental estimates used to reach the finite interface are written
as directed rational inequalities.  Two short Python ledgers check the 216
terminal cases and the final cross products using \texttt{fractions.Fraction}.
The complete route is also formalized in Lean 4
\cite{deMouraEtAl2021,Mathlib2020}.  The human proof, executable arithmetic,
and formal proof are deliberately separate evidence layers.

\section{The analytic-to-matrix transfer}

Because the new argument modifies the rank--trace step of
\cite[Sections~2--7]{Claude2026}, we spell out the complete interface used
here.  In particular, the labels, matrices, normalization, boundary loss and
fixed-window comparison are not left implicit.

For $0<\lambda\le1$ put
\begin{equation}\label{eq:fixed-lambda-constant}
 \vartheta_\lambda=\frac{\lambda}{\sqrt2},\qquad
 c_\lambda^*=\frac{\sqrt2\sin\vartheta_\lambda}
 {\cos\vartheta_\lambda+\vartheta_\lambda\sin\vartheta_\lambda},
 \qquad H_\lambda=2-(c_\lambda^*)^{-1}.
\end{equation}
Thus $H_1=H_*$.  Fix $\lambda<1$ while $T\to\infty$, and write
\begin{align*}
 \ell&=\log(T/2\pi), & L&=\lambda\ell, & h&=2\pi/L,\\
 d&=\lfloor LT/(2\pi)\rfloor,
 & \tau_k&=T+kh &&(0\le k<d).
\end{align*}
Let $w=1$ and let $\chi$ be a fixed nondecreasing $C^3$ ramp which is zero to
the left of $0$ and one to the right of $1$; in particular,
$0\le\chi\le1$.  We use the Montgomery--Taylor window
\begin{equation}\label{eq:finite-window}
 \phi_T(u)=
 \sqrt{\cos(\sqrt2\lambda u/L)}\,
 \chi\left(\frac{L/2-|u|}{w}\right),
 \qquad
 a_T=\frac1L\int_{\mathbb R}\phi_T(u)^2\,du,
\end{equation}
and the Fourier convention
$\widehat\phi_T(r)=\int\phi_T(u)e^{iru}\,du$.  On the support of the
window the cosine in \eqref{eq:finite-window} is at least $3/4$, and hence
$1/2\le a_T\le1$ for all sufficiently large $T$.

For a zero $\rho$, set $\gamma_\rho=-i(\rho-1/2)$; it is real precisely on
the critical line, and the functional equation sends $\gamma_\rho$ to
$\overline{\gamma_\rho}$.  The finite Weil matrix is the Hermitian matrix of
the zero-side form on the modulated windows
$\phi_T(u)e^{-i\tau_ku}$.  Summing over the distinct nontrivial zeros once
and writing their multiplicities explicitly, this says
\begin{equation}\label{eq:weil-entry}
 (\mathcal G_T)_{kl}=\sum_\rho m_\rho\,
 \widehat\phi_T(\gamma_\rho-\tau_k)
 \overline{\widehat\phi_T(\overline{\gamma_\rho}-\tau_l)}.
\end{equation}
Pairing $\rho$ with $1-\overline\rho$ makes
\eqref{eq:weil-entry} Hermitian.  We normalize it by
\begin{equation}\label{eq:weil-normalization}
 \widehat{\mathcal G}_T=\frac{\mathcal G_T}{a_TL^2}.
\end{equation}
For a simple critical-line zero $\rho=1/2+i\gamma$, its contribution to
$\widehat{\mathcal G}_T$ is the rank-one matrix $v_\gamma v_\gamma^*$, where
\begin{equation}\label{eq:finite-vector}
 v_\gamma(k)=\frac{\widehat\phi_T(\gamma-\tau_k)}{L\sqrt{a_T}},\qquad
 w_\gamma=\|v_\gamma\|_2^2,\qquad
 x_\gamma=\frac{v_\gamma}{\sqrt{w_\gamma}}.
\end{equation}
The last expression is used only after $w_\gamma>0$ has been established.

Define the interior simple core
\begin{align}\label{eq:simple-core}
 \mathcal S_T&=\left\{\frac12+i\gamma:
   \zeta\left(\frac12+i\gamma\right)=0,\ m_\rho=1,\right.\notag\\[-2mm]
 &\hspace{39mm}\left.T+3\le\gamma\le2T-3\right\},
 & s=s_T=|\mathcal S_T|.
\end{align}
Put $I'=(T-\sqrt T,2T+\sqrt T]$, let $s_1$ count the simple critical-line
zeros in $I'$, and set $E_{\partial}=s_1-s$.  Thus $E_{\partial}$ counts the
simple labels removed from $I'$ when the core is selected.  Also put
\begin{equation}\label{eq:dimension-D}
 D=D_T=\frac{LT}{2\pi}.
\end{equation}
The matrix $G_T=(\langle x_\gamma,x_{\gamma'}\rangle)_
{\gamma,\gamma'\in\mathcal S_T}$ is the simple-core Gram matrix; it has unit
diagonal.  The simple zeros in $I'$ but omitted from the core lie in two
outer strips of total length $2\sqrt T$ and two interior strips of length
three.  The local zero-count
bound $N(t,t+1)=O(\log(|t|+3))$ therefore gives
\begin{equation}\label{eq:boundary-loss}
 E_{\partial}(T)=O(\sqrt T\log T)=o(N(T,2T)).
\end{equation}

We next record the analytic matrix estimates in the exact normalization just
defined.

\begin{lemma}[Fixed-window matrix package]\label{lem:matrix-package}
For every fixed $0<\lambda<1$, as $T\to\infty$,
\begin{align}
 \tr\widehat{\mathcal G}_T&=N(T,2T)+o(N(T,2T)),\label{eq:trace-one}\\
 \|\widehat{\mathcal G}_T\|_F^2
   &=(c_\lambda^*)^{-1}N(T,2T)+o(N(T,2T)),\label{eq:trace-two}\\
 \frac{s}{N(T,2T)}&\ge H_\lambda-o(1),\qquad
 \frac{D}{N(T,2T)}=\lambda+o(1),\qquad
 \frac{s}{N(T,2T)}\le1.\label{eq:parent-density}
\end{align}
The normalized tail outside $I'$ has trace norm $o(1)$.
No assertion here assumes RH or a zero-density estimate.
\end{lemma}

\begin{proof}
Poisson summation at spacing $h=2\pi/L$ gives, without aliasing,
\[
 \sum_{k\in\mathbb Z}
 \widehat\phi_T(\gamma-\tau_k)
 \overline{\widehat\phi_T(\gamma'-\tau_k)}
 =L\,\widehat{\phi_T^2}(\gamma-\gamma').
\]
The diagonal case is $a_TL^2$.  Applying Weil's explicit formula to the
finite family and using the Montgomery--Vaughan inequality for the prime
powers $n\le e^L=(T/2\pi)^\lambda$ gives
\eqref{eq:trace-one}.  The same calculation for the squared matrix gives
\[
 \|\widehat{\mathcal G}_T\|_F^2
 =\frac{b_T+\lambda^2J_T}{\lambda a_T^2}\,N(T,2T)+o(N(T,2T)),
\]
where $b_T=L^{-1}\int\phi_T^4$ and
$J_T=2L^{-3}\int_0^L(\phi_T^2\star\phi_T^2)(y)y\,dy$.
The fixed-width ramps have total mass $O(1/L)$, so
$(a_T,b_T,J_T)$ converge to the scale-free moments of
$v_\lambda(s)=\cos(\sqrt2\lambda s)$.  The elementary Euler--Lagrange
calculation for this profile gives
$\lambda a_T^2/(b_T+\lambda^2J_T)\to c_\lambda^*$, proving
\eqref{eq:trace-two}.  These are the trace calculations in
\cite[Sections~2, 5 and 7.1]{Claude2026}; the preceding display states all
normalizations needed here.

On the zero side, write
$\widehat{\mathcal G}_T=\widehat{\mathcal A}_T+\widehat{\mathcal E}_T$,
where $\widehat{\mathcal A}_T$ is the contribution of the zeros with real
ordinate in $I'$ and $\widehat{\mathcal E}_T$ is the tail.  Each
critical-line point contributes a positive rank-one
form.  Each off-line functional-equation pair contributes the pullback of a
$2\times2$ form of signature $(1,1)$, so its positive index is at most one.
For Hermitian $P\succeq0$ and Hermitian $Q$, with
$\rank P\le r$ and $n_+(Q)\le b$, von Neumann's trace inequality yields
\begin{equation}\label{eq:rank-trace-basic}
 r\ge2\tr P+4\tr Q-4b-\|P+Q\|_F^2.
\end{equation}
Indeed, writing $Q=Q_+-Q_-$ and arranging the eigenvalues of $P,Q_+,Q_-$
in decreasing order, von Neumann bounds the only unfavorable cross term by
$\sum_i p_i n_i$; the scalar inequalities
$(p_i-n_i)^2\ge2(p_i-n_i)-1$ and $q_j^2\ge4q_j-4$ then sum to
\eqref{eq:rank-trace-basic}.  Repeated eigenvalues follow by continuity.
Applying \eqref{eq:rank-trace-basic} to $\widehat{\mathcal A}_T$, with the
simple critical-line points on the rank side and the multiple on-line points
and off-line pairs on the index side, gives
\[
 s_1
 \ge4\tr\widehat{\mathcal A}_T-
 \|\widehat{\mathcal A}_T\|_F^2-2N(I').
\]
Proposition~4.2 of \cite{Claude2026} gives
$\|\widehat{\mathcal E}_T\|_1=O_\lambda(T^{\lambda/2-1})$.
Consequently $|\tr\widehat{\mathcal E}_T|=o(N)$ and
\[
 \left|\|\widehat{\mathcal A}_T\|_F^2-
 \|\widehat{\mathcal G}_T\|_F^2\right|
 \le \|\widehat{\mathcal E}_T\|_1
 \left(2\|\widehat{\mathcal G}_T\|_F+
 \|\widehat{\mathcal E}_T\|_1\right)=o(N).
\]
Also $N(I')=N(T,2T)+O(\sqrt T\log T)$.  Equations
\eqref{eq:trace-one}--\eqref{eq:trace-two}, followed by
$s=s_1-E_{\partial}$ and \eqref{eq:boundary-loss}, therefore give the first
inequality in \eqref{eq:parent-density}.  Whenever $H_\lambda>0$, this first,
defect-free pass also proves that the core is nonempty for all large $T$.
The second relation in \eqref{eq:parent-density} follows from
$D=LT/(2\pi)$ and the Riemann--von Mangoldt formula; the last is exact because
$\mathcal S_T\subset(T,2T]$.  The displayed tail estimate is exactly the
term-by-term consequence of two integrations by parts in $\widehat\phi_T$
and the unit-window zero count; see \cite[Proposition~4.2]{Claude2026}.
\end{proof}

The four-point improvement enters through the following finite-height
version of the same zero-side calculation.  Here $\delta_\lambda$ is a
certified \emph{lower bound} for the finite/full four-point energy, not the
energy itself.

\begin{lemma}[Refined unit-core rank--trace inequality]
\label{lem:refined-rank-trace}
Let $A=P+Q$ be Hermitian, where
$P=\sum_{j=1}^s x_jx_j^*$, where $s\ge1$ and $\|x_j\|=1$, and suppose
$n_+(Q)\le b$.  If $G=(\langle x_i,x_j\rangle)_{i,j\le s}$ and
$V=\|G-I_s\|_1/2$, then
\begin{equation}\label{eq:refined-unit-rank-trace}
 \|A\|_F^2\ge
 2\tr P-s+4\tr Q-4b+\frac{V^2}{s}.
\end{equation}
\end{lemma}

\begin{proof}
Let $d$ be the ambient dimension and, for a positive semidefinite matrix
$R$ on $\mathbb C^m$, put
\[
 \mathcal D_m(R)=\sum_{i=1}^m(1-\lambda_i(R))_+^2.
\]
Splitting $Q=Q_+-Q_-$, von Neumann's trace inequality and
\[
 (p-n)^2\ge2p-1-4n+(1-p)_+^2,
 \qquad q^2\ge4q-4
\]
give, after summing the ambient $d$ eigenvalue contributions,
\[
 \|P+Q\|_F^2\ge
 2\tr P-d+4\tr Q-4b+\mathcal D_d(P).
\]
This is where the change from the ambient spectrum to the Gram spectrum must
be recorded.  If $W$ has columns $x_1,\ldots,x_s$, then
$P=WW^*$ and $G=W^*W$.  If their common nonzero eigenvalues are
$\mu_1,\ldots,\mu_r$, their zero multiplicities are $d-r$ and $s-r$,
respectively.  Since $(1-0)_+^2=1$, it follows exactly that
\[
 \mathcal D_d(P)-d
 =\sum_{i=1}^r(1-\mu_i)_+^2-r
 =\mathcal D_s(G)-s.
\]
Substitution therefore changes the preceding ambient bound into
\[
 \|P+Q\|_F^2\ge
 2\tr P-s+4\tr Q-4b+\mathcal D_s(G).
\]
Let $p_1,\ldots,p_s\ge0$ be the eigenvalues of $G$, including its zero
eigenvalues.  Then $\mathcal D_s(G)=\sum_j(1-p_j)_+^2$.
The unit diagonal gives $\sum_jp_j=s$, hence the total mass below one is
\[
 \sum_j(1-p_j)_+=\frac12\sum_j|p_j-1|
 =\frac12\|G-I_s\|_1=V.
\]
Cauchy--Schwarz now gives
$\sum_j(1-p_j)_+^2\ge V^2/s$, proving
\eqref{eq:refined-unit-rank-trace}.  This derivation does not require $P$
and $Q$ to commute.
\end{proof}

\begin{proposition}[Four-point transfer]\label{prop:transfer}
Fix $\lambda$ with $H_\lambda>0$ and a constant $\delta_\lambda\ge0$.
For each sufficiently large $T$ with $s>0$, suppose that every packed
width-six four-tuple of the core at that height has six-edge energy at least
$2\delta_\lambda$.  Then, at each height for which this local hypothesis
holds,
\begin{equation}\label{eq:four-defect}
 N_0^{\mathrm s}(T,2T)\ge
 4\tr\widehat{\mathcal G}_T-
 \|\widehat{\mathcal G}_T\|_F^2-2N(T,2T)
 +\frac{\delta_\lambda}{8s}
   \pos{s-\frac D2-3}^{2}-o(N(T,2T)).
\end{equation}
If the same $\delta_\lambda$ satisfies the local hypothesis for every
$T\ge T_1(\lambda,\delta_\lambda)$, then
\begin{equation}\label{eq:fixed-lambda-form}
 \liminf_{T\to\infty}\frac{N_0^{\mathrm{s}}(T,2T)}{N(T,2T)}
 \ge H_\lambda+
 \frac{\delta_\lambda}{8}
 \pos{H_\lambda-\frac\lambda2}^{2}.
\end{equation}
\end{proposition}

\begin{proof}
Partition the scaled ordinates $L\gamma$ into the half-open intervals
$I_j=[L(T+3)+12\pi j,L(T+3)+12\pi(j+1))$, $j\ge0$.  For every core
ordinate, $0\le L\gamma-L(T+3)\le LT$; hence the occupied indices satisfy
$0\le j\le\lfloor D/6\rfloor$.  There are therefore at most $D/6+1$ bins.
Taking consecutive groups of four in each bin loses at most three
labels per occupied bin and gives
\[
 t\ge\frac14\pos{s-D/2-3}.
\]
For the unit core vectors, let $P_0=\sum_{\gamma\in\mathcal S_T}
x_\gamma x_\gamma^*$.  The actual simple-core part is
$P_w=\sum w_\gamma x_\gamma x_\gamma^*$ with $0<w_\gamma\le1$.
Let $s_2$ be the number of distinct multiple critical-line zeros in $I'$ and
let $p$ count the off-line functional-equation pairs there.  Before
unitization, the complement
$Q_w=\widehat{\mathcal A}_T-P_w$ has positive index at most
\[
 b=E_{\partial}+s_2+p.
\]
This is the direct sum/pullback count: the excluded simple labels and
multiple on-line labels give positive one-dimensional forms, and each
off-line pair gives one positive direction.  Since
$P_0-P_w\succeq0$, the unitized complement
\[
 Q_0=\widehat{\mathcal A}_T-P_0=Q_w-(P_0-P_w)
\]
satisfies $n_+(Q_0)\le n_+(Q_w)\le b$ by the min--max principle.
Lemma~\ref{lem:refined-rank-trace} therefore applies.  Its spectral term is
$V^2/s$, where
$V=\|G_T-I_s\|_1/2$.  The four-block calculation in
Section~\ref{sec:four-point} gives
$V\ge t\sqrt{2\delta_\lambda}$ and hence
\[
 \frac{V^2}{s}\ge\frac{2t^2\delta_\lambda}{s}
 \ge\frac{\delta_\lambda}{8s}\pos{s-D/2-3}^2.
\]
Because $\tr P_0=s$ and $\tr Q_0=\tr\widehat{\mathcal A}_T-s$,
Lemma~\ref{lem:refined-rank-trace} gives
\[
 \|\widehat{\mathcal A}_T\|_F^2
 \ge4\tr\widehat{\mathcal A}_T-3s-4b+\frac{V^2}{s}.
\]
The multiplicity count
$N(I')\ge s+E_{\partial}+2s_2+2p$ implies
$4b\le2N(I')-2s+2E_{\partial}$.  Rearranging, and using
$N_0^{\mathrm s}(T,2T)\ge s$, yields
\[
 N_0^{\mathrm s}(T,2T)\ge
 4\tr\widehat{\mathcal A}_T-
 \|\widehat{\mathcal A}_T\|_F^2-2N(I')-2E_{\partial}
 +\frac{V^2}{s}.
\]
The tail and boundary estimates in Lemma~\ref{lem:matrix-package} replace
$\widehat{\mathcal A}_T$ by $\widehat{\mathcal G}_T$ and $N(I')$ by
$N(T,2T)$; they absorb $E_{\partial}$ at total cost $o(N)$.
Together with the displayed lower bound for $V^2/s$, this proves
\eqref{eq:four-defect}.  Divide it by $N=N(T,2T)$, use
Lemma~\ref{lem:matrix-package}, and note that $s/N\le1$ makes the reciprocal
factor at least one.  Under the eventual hypothesis in the statement, the
positive-part argument tends from below to
$\pos{H_\lambda-\lambda/2}$ along every sufficiently large height, proving
\eqref{eq:fixed-lambda-form}.
\end{proof}

We finally make explicit how an endpoint energy bound becomes a legitimate
fixed-window input.  Define
\begin{equation}\label{eq:fixed-kernel}
 K_\lambda(x)=\int_{-1/2}^{1/2}
   \cos(\sqrt2\lambda s)\cos(xs)\,ds,\qquad
 R_\lambda(x)=\frac{K_\lambda(x)}{K_\lambda(0)}.
\end{equation}
For the vectors \eqref{eq:finite-vector}, Poisson summation identifies the
full-grid correlation with the Fourier transform of $\phi_T^2$.  The omitted
grid tail has normalized size
\begin{equation}\label{eq:grid-tail}
 0\le q_T\le
 \frac{4c_\chi^2}{81L}+\frac{113c_\chi^2}{648L^2}=o(1),
\end{equation}
where $c_\chi$ depends only on the fixed ramp.  Finite Cauchy--Schwarz gives
$w_\gamma\in[1-q_T,1]$, and scalar normalization gives a uniform correlation
error at most $2q_T$.

The $L^1$ distance between $\phi_T^2$ and the sharp profile
$\cos(\sqrt2\lambda u/L)1_{[-L/2,L/2]}$ is at most $2w$.  The numerator error
is therefore at most $2w/L$, the zero-frequency mass error at most $4w/L$,
and the actual mass is at least $1/2$ when $16w\le L$.  Hence the normalized
full-grid kernel differs from $R_\lambda$ by at most $12w/L$.  Moreover
\[
 |\cos(\sqrt2\lambda s)-\cos(\sqrt2s)|\le1-\lambda
 \quad(|s|\le1/2),
\]
while $K_\lambda(0)\ge3/4$.  Comparing numerator and denominator in
\eqref{eq:fixed-kernel} gives, uniformly in $x$,
\begin{equation}\label{eq:kernel-comparison}
 |R_\lambda(x)-R_1(x)|\le3(1-\lambda).
\end{equation}
Combining the three comparisons, every finite normalized Gram correlation is
within
\begin{equation}\label{eq:correlation-error}
 \epsilon_T=\frac{12w}{L}+3(1-\lambda)+2q_T
 \longrightarrow3(1-\lambda)
\end{equation}
of the endpoint correlation.  Since $|u|,|v|\le1$ and $|u-v|\le\epsilon_T$
imply $v^2\ge u^2-2\epsilon_T$, a three-edge energy loses at most
$6\epsilon_T$.  Therefore the actual finite-height four-point defect
parameter tends to
\begin{equation}\label{eq:eighteen-loss}
 \delta_*-6\cdot3(1-\lambda)=\delta_*-18(1-\lambda).
\end{equation}
This proves, rather than assumes, the fixed-window comparison used in
Section~\ref{sec:endpoint}.

\section{The endpoint kernel}

The endpoint profile is
\[
 v(s)=\cos(\sqrt2s),\qquad -\frac12\le s\le\frac12.
\]
Define the real Fourier kernel and its normalization by
\begin{equation}\label{eq:kernel-def}
 K(x)=\int_{-1/2}^{1/2}\cos(\sqrt2s)e^{ixs}\,ds,
 \qquad R(x)=\frac{K(x)}{K(0)}.
\end{equation}
The imaginary part vanishes by parity.  Put
\begin{equation}\label{eq:kappa}
 \kappa=\sqrt2\tan(1/\sqrt2).
\end{equation}
Away from the removable points $x^2=2$, product-to-sum gives
\begin{equation}\label{eq:kernel-closed}
 K(x)=\frac{2\cos(1/\sqrt2)}{x^2-2}
 \left(x\sin\frac x2-\kappa\cos\frac x2\right).
\end{equation}
Moreover $0<K(0)\le1$, and hence $|R(x)|\ge|K(x)|$.

For $n\ge1$, let
\begin{equation}\label{eq:roots}
 z_n=2\pi n+\eps_n
\end{equation}
be the positive root in the $n$th positive half-period.  Equivalently,
\begin{equation}\label{eq:root-equation}
 (2\pi n+\eps_n)\tan(\eps_n/2)=\kappa.
\end{equation}
The inherited root package gives
\begin{equation}\label{eq:offset-basic}
 0<\eps_{n+1}<\eps_n<\frac25,
 \qquad \eps_3>\frac18.
\end{equation}

\section{Exact root neighborhoods}

Set
\begin{equation}\label{eq:rho-tau}
 \rho=\frac18,
 \qquad \tau=\frac{1815107}{989072150}.
\end{equation}
For $1\le n\le6$ define
\[
 g_n(e)=(2\pi n+e)\sin(e/2)-\kappa\cos(e/2).
\]
Its derivative is
\begin{equation}\label{eq:gn-derivative}
 g_n'(e)=\left(1+\frac\kappa2\right)\sin\frac e2
       +\frac{2\pi n+e}{2}\cos\frac e2.
\end{equation}

We use the rational bounds
\begin{equation}\label{eq:elementary-bounds}
 \frac65<\kappa\le\frac{49}{40},\qquad
 3<\pi<\frac{22}{7},\qquad
 2\cos(1/\sqrt2)\ge\frac32.
\end{equation}
On $-1/8\le e\le21/40$, one has
\[
 \cos(e/2)\ge C:=\frac{12359}{12800}
\]
and therefore
\begin{equation}\label{eq:Ln}
 g_n'(e)\ge L_n:=
 \left(3n-\frac1{16}\right)\frac{12359}{12800}
 -\frac{129}{1280}>0.
\end{equation}
Integration from the root, followed by \eqref{eq:kernel-closed}, gives the
following exact local slopes.

\begin{lemma}[Linear root neighborhoods]\label{lem:slopes}
If $1\le n\le6$ and $|x-z_n|\le1/8$, then
\begin{equation}\label{eq:local-slope}
 |R(x)|\ge \frac{|x-z_n|}{d_n},
\end{equation}
where
\begin{equation}\label{eq:dn-table}
\begin{aligned}
d_1&=\frac{930982144}{82354251},&
d_2&=\frac{3442403584}{169559355},&
d_3&=\frac{2513265408}{85588153},\\
d_4&=\frac{13223160064}{343969563},&
d_5&=\frac{20492495104}{431174667},&
d_6&=\frac{9782600448}{172793257}.
\end{aligned}
\end{equation}
\end{lemma}

\begin{proof}
If $x=2\pi n+e$, the segment joining $e$ to $\eps_n$ stays in the strip on
which \eqref{eq:Ln} holds.  Thus
$|g_n(e)|\ge L_n|e-\eps_n|$.  On the same strip,
$x<44n/7+21/40$.  Equations \eqref{eq:kernel-closed} and
\eqref{eq:elementary-bounds} imply
\[
 |R(x)|\ge|K(x)|\ge
 \frac{3L_n}{2(44n/7+21/40)^2}|e-\eps_n|.
\]
Exact reduction of the six rational reciprocals gives
\eqref{eq:dn-table}.
\end{proof}

Two offset refinements will be needed.  The elementary tangent bound
$\tan u\le u/(1-u^2/2)$ gives
\[
 \left(4\pi+\frac{11}{60}\right)\tan\frac{11}{120}
 <\frac{235708}{200753}<\frac65<\kappa,
\]
so monotonicity of \eqref{eq:root-equation} yields
\begin{equation}\label{eq:eps2}
 \eps_2>\frac{11}{60}.
\end{equation}
Conversely,
\[
 \left(12\pi+\frac1{15}\right)\tan\frac1{30}
 >\frac{2831}{2250}>\frac{49}{40}\ge\kappa,
\]
and hence
\begin{equation}\label{eq:eps6}
 \eps_6<\frac1{15}.
\end{equation}

\section{Six-root localization}

\begin{proposition}[Localization]\label{prop:localization}
For $0\le x\le12\pi$,
\begin{equation}\label{eq:localization}
 |R(x)|<\tau
 \quad\Longrightarrow\quad
 \min_{1\le n\le6}|x-z_n|<\frac18.
\end{equation}
\end{proposition}

\begin{proof}
It is enough to prove the contrapositive with $K$, because
$|K(x)|\le|R(x)|$.  The interval is split at multiples of $\pi$.

On $[0,\pi]$, the inherited initial estimate gives
$K(x)\ge165/512>\tau$.  On each positive half-period
$[2\pi n,(2n+1)\pi]$, $1\le n\le5$, if the point is outside the radius
$1/8$ neighborhood then monotonicity and Lemma~\ref{lem:slopes} give
\begin{equation}\label{eq:positive-half}
 |K(x)|\ge
 \frac32\frac{L_n/8}{(22(2n+1)/7)^2}>\tau.
\end{equation}
The weakest comparison is $n=5$, for which the positive slack is
\[
 \frac{3563427}{1585971200}-\tau
 =\frac{106740326833}{259279329689600}>0.
\]

On the five zero-free half-periods $[(2m-1)\pi,2m\pi]$, $1\le m\le5$,
the two numerator terms in \eqref{eq:kernel-closed} have the same sign.  The
refined bound $\kappa\ge592688/490449$ gives
\begin{equation}\label{eq:zero-free}
 |K(x)|\ge
 \frac32\frac{592688/490449}{(44m/7)^2}\ge\tau.
\end{equation}
At $m=5$ equality holds.  This is sufficient because the premise in
\eqref{eq:localization} is strict.

It remains to treat $[11\pi,12\pi]$.  Write $x=12\pi+e$ with
$-\pi\le e\le0$.  If $e\le-1/8$, then
$\sin(-e/2)\ge\sin(1/16)>1/17$, and
\[
 |K(x)|\ge\frac32\frac{33/17}{(264/7)^2}
 =\frac{49}{23936}>\tau.
\]
If $-1/8<e\le\eps_6-1/8$, integrate \eqref{eq:Ln} over the final length
$1/8$ to obtain
\[
 |K(x)|\ge\frac32\frac{L_6/8}{(264/7)^2}
 =\frac{172793257}{76126617600}>\tau.
\]
Together with the already localized neighborhood $|x-z_6|<1/8$, these
estimates exhaust the points of $[0,12\pi]$ that must be considered.
\end{proof}

\section{The wider three-coordinate energy}

Define
\begin{equation}\label{eq:delta16-def}
 \delta_{1,6}:=
 \min_{\substack{a,b\ge0\\a+b\le12\pi}}
 \bigl(R(a)^2+R(b)^2+R(a+b)^2\bigr).
\end{equation}

\begin{theorem}[Exact local energy]\label{thm:local-energy}
One has
\begin{equation}\label{eq:delta-star}
 \delta_{1,6}\ge\delta_*:=
 \frac{2810232522752}{834437357315980975}.
\end{equation}
\end{theorem}

\begin{proof}
Suppose the energy in \eqref{eq:delta16-def} were smaller than $\delta_*$.
Exact subtraction shows $\delta_*<\tau^2$, so each of
$|R(a)|,|R(b)|,|R(a+b)|$ is smaller than $\tau$.  By
Proposition~\ref{prop:localization}, there are $i,j,k\in\{1,\ldots,6\}$
such that
\[
 |a-z_i|<\frac18,\quad |b-z_j|<\frac18,
 \quad |a+b-z_k|<\frac18.
\]
Consequently $|z_i+z_j-z_k|<3/8$.  If $i+j\ne k$, the nonzero integral
multiple of $2\pi$ dominates the three offsets and gives
$|z_i+z_j-z_k|>2\pi-4/5>3/8$, a contradiction.  Hence $i+j=k$.

Up to swapping $i$ and $j$, there are eight additive triples other than
$(3,3,6)$.  For them, monotonicity of the offsets and \eqref{eq:eps2} give
\[
 \eps_i+\eps_j-\eps_k>\frac{11}{60}.
\]
For $(3,3,6)$, equations \eqref{eq:offset-basic} and \eqref{eq:eps6} give
the same floor:
$2\eps_3-\eps_6>1/4-1/15=11/60$.

Put $u=a-z_i$, $v=b-z_j$ and $w=a+b-z_k$.  Lemma~\ref{lem:slopes} and
weighted Cauchy--Schwarz yield
\begin{align}\label{eq:weighted-cauchy}
 R(a)^2+R(b)^2+R(a+b)^2
 &\ge \frac{u^2}{d_i^2}+\frac{v^2}{d_j^2}+\frac{w^2}{d_k^2},\\
 (z_i+z_j-z_k)^2
 &\le(d_i^2+d_j^2+d_k^2)
 \left(\frac{u^2}{d_i^2}+\frac{v^2}{d_j^2}
 +\frac{w^2}{d_k^2}\right).\notag
\end{align}
For completeness, set $\sigma_{ijk}=11/60$ when $i+j=k$ and
$\sigma_{ijk}=1$ otherwise.  The 216 exact budgets are
\[
 B_{ijk}=\frac{\sigma_{ijk}^2}{d_i^2+d_j^2+d_k^2}.
\]
Their minimum is attained at $(1,5,6)$ and $(5,1,6)$ and equals
\begin{equation}\label{eq:budget-min}
 B_*=
 \frac{9755468468368320337833317374859207102505001}
 {1623000197669009639771565694101668583197612441600}.
\end{equation}
The exact ledger verifies $\delta_*<\tau^2<B_*$.  In particular,
\[
 \tau^2-\delta_*
 =\frac{35399897577222712511}
 {32651991661090957509227042221277500}>0.
\]
Thus \eqref{eq:weighted-cauchy} contradicts the assumed energy, proving the
theorem.
\end{proof}

\section{Four-point packing and spectral defect}\label{sec:four-point}

The reason for working on the wider interval $[0,12\pi]$ is that it supports
four-point rather than three-point packing.  The following elementary
identity is the source of the improvement.

\begin{lemma}[Four points count six edges twice]\label{lem:four-energy}
Let $x_1,x_2,x_3,x_4$ lie in a normalized interval of width $12\pi$.
If each three-coordinate energy is at least $\delta$, then
\begin{equation}\label{eq:four-energy}
 2\delta\le
 \sum_{1\le i<j\le4}|R(x_i-x_j)|^2.
\end{equation}
\end{lemma}

\begin{proof}
Apply the three-coordinate inequality to each of the four triples obtained by
deleting one point.  Every one of the six unordered pairs belongs to exactly
two of these triples.  Summing the four inequalities and dividing by two
gives \eqref{eq:four-energy}.
\end{proof}

The spectral conversion is also sharp for this block size.

\begin{lemma}[Trace-zero four-block]\label{lem:trace-zero}
If $B$ is Hermitian and $\tr B=0$, then
\begin{equation}\label{eq:trace-zero}
 \|B\|_1^2\ge2\|B\|_F^2.
\end{equation}
\end{lemma}

\begin{proof}
Let $p$ be the sum of the positive eigenvalues.  Trace zero implies that the
absolute sum of the negative eigenvalues is also $p$, so $\|B\|_1=2p$.
The sum of squares on either sign is at most the square of the corresponding
sum.  Hence $\|B\|_F^2\le2p^2=\|B\|_1^2/2$.
\end{proof}

We spell out the packing count.  If $s$ simple core labels are distributed
among width-six bins and the parent dimensional budget is $D$, the number of
occupied bins is at most $D/6+1$.  Partition each bin into disjoint
four-tuples, leaving at most three labels in that bin.  If $t$ four-tuples are
obtained, then
\[
 4t\ge s-3\left(\frac D6+1\right)
 =s-\frac D2-3.
\]
Thus
\begin{equation}\label{eq:t-count}
 t\ge\frac14\pos{s-\frac D2-3}.
\end{equation}

Let $G$ be the simple-core Gram matrix and $H=G-I_s$.  Thus $H$ is Hermitian
and trace zero.  On each selected four-tuple, Lemma~\ref{lem:four-energy}
and the zero diagonal of $H$ give
\[
 \|H_{\mathrm{block}}\|_F^2
 =2\sum_{i<j}|G_{ij}|^2\ge4\delta_\lambda.
\]
Lemma~\ref{lem:trace-zero} therefore gives
$\|H_{\mathrm{block}}\|_1\ge2\sqrt{2\delta_\lambda}$.
Pinching onto disjoint blocks is an average of unitary conjugations, so it is
contractive for the trace norm.  If
$V=\frac12\|H\|_1$ is the negative spectral mass, then
\[
 V\ge t\sqrt{2\delta_\lambda}.
\]
The refined rank--trace lemma of the parent framework contributes at least
$V^2/s$ to its defect.  Using \eqref{eq:t-count},
\[
 \frac{V^2}{s}\ge\frac{2t^2\delta_\lambda}{s}
 \ge\frac{\delta_\lambda}{8s}
 \pos{s-\frac D2-3}^{2},
\]
which is \eqref{eq:four-defect}.  This is exactly the improvement over the
earlier triple coefficient.  The full matrix statement, including degenerate
blocks and the positive-index comparison, is formalized in
\texttt{FourBlockDefect.lean}.

\section{Endpoint passage}\label{sec:endpoint}

At a fixed $\lambda<1$, the normalized finite-window kernel converges to its
scale-free kernel, while the scale-free kernel converges uniformly to the
endpoint kernel as $\lambda\to1^-$.  The six correlations in a four-tuple
produce a total comparison loss at most $18(1-\lambda)$.  Consequently, for
every fixed $\lambda$ sufficiently close to one and every constant satisfying
\begin{equation}\label{eq:fixed-delta-choice}
 0\le\delta_\lambda<\delta_*-18(1-\lambda),
\end{equation}
$\delta_\lambda$ is a valid lower-bound input to
Proposition~\ref{prop:transfer}.  The strict inequality leaves room for the
finite-window approximation.

The order of limits is important.  Put
\[
 h_n=\frac1{n+1},\qquad \lambda_n=1-h_n,qquad
 \delta_n=\delta_*-19h_n.
\]
For all sufficiently large $n$, $\delta_n\ge0$ and
$\delta_n<\delta_*-18(1-\lambda_n)$, so
\eqref{eq:fixed-delta-choice} holds.  Moreover,
$H_{\lambda_n}-\lambda_n/2\to H_*-1/2>0$, so the positive part in
\eqref{eq:fixed-lambda-form} is eventually inactive.  Given the outer error $\eta>0$ in
Theorem~\ref{thm:main}, convergence of the resulting coefficients lets us
choose one sufficiently large index $n=n(\eta)$ whose coefficient is within
$\eta/2$ of the endpoint value.  We then fix this $n$ and invoke the
fixed-$\lambda_n$ theorem with its remaining error budget $\eta/2$, obtaining
$T_0=T_0(n,\eta)$.  The asserted inequality holds for every $T\ge T_0$.
Thus the quantifier order is ``choose $n$ from $\eta$, then choose $T_0$'';
there is no diagonal choice $\lambda=\lambda(T)$ and no exchange of limits.

Since $H_{\lambda_n}\to H_*$ and $\delta_n\to\delta_*$, the endpoint gain is
\begin{equation}\label{eq:endpoint-gain}
 \frac{\delta_*}{8}\left(H_*-\frac12\right)^2.
\end{equation}
The exact endpoint estimate
\begin{equation}\label{eq:endpoint-gap}
 H_*-\frac12\ge\frac{102239}{592688}
\end{equation}
and rational reduction give
\begin{equation}\label{eq:eta-exact}
 \frac{\delta_*}{8}
 \left(\frac{102239}{592688}\right)^2
 =\frac1{79\,828\,975}.
\end{equation}
Equations \eqref{eq:endpoint-gain}--\eqref{eq:eta-exact}, together with the
quantified fixed-$\lambda$ transfer, prove Theorem~\ref{thm:main}.

\section{The ceiling of the frozen finite interface}

The preceding proof deliberately uses a finite rational interface.  Its
limiting obstruction can therefore be identified exactly.  Let
\begin{equation}\label{eq:activation}
 A_*=\tau^2=\frac{3294613421449}{978263717905622500}.
\end{equation}
Equation \eqref{eq:budget-min} satisfies $A_*<B_*$.  Hence the weighted
Cauchy budgets are not active: the strict localization implication
$|R|<\tau$ forces the admissible local constant to satisfy $\delta<A_*$.
The supremal endpoint gain within this interface is therefore
\begin{equation}\label{eq:Jstar}
 J_*:=\frac{A_*}{8}
 \left(\frac{102239}{592688}\right)^2
 =\frac{25097204303521}{2003484094270714880000}.
\end{equation}
Exact comparison yields
\begin{equation}\label{eq:unit-ceiling}
 \frac1{79\,828\,975}<J_*\le\frac1{79\,828\,974}.
\end{equation}

\begin{proposition}[Fixed-interface ceiling]\label{prop:ceiling}
The unit fraction in Theorem~\ref{thm:main} is the largest reciprocal integer
strictly below the endpoint supremum of the frozen radius-$1/8$, six-root,
weighted-Cauchy, four-point interface.  In particular, the same interface
cannot certify the increment $1/79\,828\,974$.
\end{proposition}

This is not an upper bound for the true simple-zero proportion, for all
four-point arguments, or even for the same kernel with sharper localization.
An improvement can come from a stronger zero-free estimate, a different
radius or block structure, a better endpoint-gap input, or a method that does
not localize each correlation separately.

\section{Formal verification and reproducibility}

The Lean development is hosted at
\url{https://github.com/aster2024/zeta-23-lean}.  The publication branch is
\begin{center}
\small\path{codex/wider-strict-improvement-79828975}.
\end{center}
The pinned source commit is
\begin{center}
\small\path{d368cc49d689576c81d7e4df0c135057d2462a27}
\end{center}
It was compiled with Lean $4.33.0$-rc2.  The pinned Mathlib commit is
\begin{center}
\small\path{51e6992efd06126df61a496bebf8f49482a4e129}
\end{center}
Public GitHub Actions run
\href{https://github.com/aster2024/zeta-23-lean/actions/runs/31637227865}{31637227865}
compiled the wider targets and the repository root.  The axiom comparator
found the same axiom union as the baseline and no occurrence of
\texttt{sorryAx} or \texttt{unsafeAx}.

The public theorem is:
\begin{verbatim}
theorem zeta_wider_strict_simple_endpoint_rational :
    forall outer > 0, exists T0 : R, forall T >= T0,
      (ThmD.HD 1 + (1 : R) / 79828975 - outer) *
        (Ncount T (2 * T) : R) <= N0simple T (2 * T)
\end{verbatim}
The actual source uses Lean's Unicode notation; the ASCII display above is
only typographical.  The theorem has no caller-supplied hypotheses.  The
repository keeps its analytic interface visible through a structure named
\texttt{PaperInputs}, then instantiates and discharges those fields before the
public theorem.  This organization is useful for auditing, but kernel
checking is not a substitute for independent expert review of the
mathematical modeling choices.

Table~\ref{tab:lean-map} maps the principal human claims to formal files.
\begin{table}[ht]
\centering
\caption{Main proof-to-Lean map.}\label{tab:lean-map}
\begin{tabular}{@{}ll@{}}
\toprule
Claim & Lean source \\
\midrule
Kernel and closed form & \texttt{EndpointKernel.lean} \\
Six-root localization & \texttt{WiderRootLocalization.lean} \\
216 energy budgets & \texttt{WiderFourPointEnergy.lean} \\
Four-block trace defect & \texttt{FourBlockDefect.lean} \\
Fixed-$\lambda$ theorem & \texttt{ZetaWiderStrictFixedLambda.lean} \\
Endpoint theorem & \texttt{ZetaWiderEndpointPassage.lean} \\
Interface ceiling & \texttt{WiderFixedInterfaceCeiling.lean} \\
\bottomrule
\end{tabular}
\end{table}

The accompanying archive contains two exact rational scripts, their frozen
outputs, and the five Lean source snapshots parsed by the first script.  They
can be replayed with a standard Python~3 interpreter:
\begin{center}\small
\texttt{python3 wider\_four\_point\_rational\_ledger.py}\\
\texttt{\hspace*{12mm}--verify-frozen-output}\\
\texttt{python3 wider\_fixed\_interface\_ceiling\_ledger.py}\\
\texttt{\hspace*{12mm}--verify-frozen-output}
\end{center}
Each script exits nonzero if a certified comparison fails.  The first reports
\texttt{root\_triples\_checked=216}; the second records the two minimizers
$(1,5,6)$ and $(5,1,6)$ and both unit-fraction cross slacks.

\section{Scope, limitations, and provenance}

The theorem is unconditional in the same mathematical sense as the parent
trace-and-pair-correlation theorem \cite{Claude2026}: it assumes neither the
Riemann hypothesis nor a zero-density hypothesis.  It is a strict improvement
of that unconditional $H_*$ bound, not of the best conditional bounds that
use more pair-correlation information.  The improvement is also not claimed
as externally established until the parent theorem and this companion
argument receive independent human expert review.

Automated systems assisted with exploration, exact-arithmetic scripting,
formalization, drafting, and adversarial proof checking.  The named human
author selected the claims, preserved the evidence and certification
boundaries, and takes responsibility for the submitted manuscript.  No AI
system is listed as an author.  The parent theorem is cited under the
authorship shown on its public manuscript.

\section*{Data and code availability}

The Lean source and workflow are available in the public repository cited
above.  The exact ledgers, their source snapshots, frozen outputs, and build
receipt are included in the companion evidence archive.  The release package
records SHA-256 hashes for its three distributed artifacts.  No proprietary
data, company MaaS key, or nonpublic model output is required to replay the
certificates.

\section*{Acknowledgments}

This work depends on the analytic framework and formal source of
\cite{Claude2026}, and on the number-theoretic literature cited there.  The
author thanks the maintainers of Lean and Mathlib for the proof environment.

\bibliographystyle{amsplain}
\bibliography{references}

\appendix

\section{Exact constants}

For convenient comparison, Table~\ref{tab:constants} lists the constants
that enter the finite interface.  Every row is present verbatim in one of the
two executable ledgers.

\begin{table}[ht]
\centering
\caption{Frozen exact constants.}\label{tab:constants}
\begin{tabular}{@{}ll@{}}
\toprule
Quantity & Exact value \\
\midrule
Localization radius $\rho$ & $1/8$ \\
Threshold $\tau$ & $1815107/989072150$ \\
Offset floor & $11/60$ \\
Local energy $\delta_*$ & $2810232522752/834437357315980975$ \\
Endpoint gap lower bound & $102239/592688$ \\
Public increment & $1/79828975$ \\
Next excluded unit & $1/79828974$ \\
\bottomrule
\end{tabular}
\end{table}

\section{Verification identities}
\enlargethispage{5\baselineskip}

The two terminal identities used in the theorem and the method ceiling are
\[
 \frac{2810232522752}{834437357315980975}\cdot\frac18
 \left(\frac{102239}{592688}\right)^2
 =\frac1{79828975}
\]
and
\[
 \frac1{79828975}<
 \frac{25097204303521}{2003484094270714880000}
 \le\frac1{79828974}.
\]
They are checked by integer cross multiplication in both Lean and Python.

\end{document}
